% GENERATED FILE. Edit canonical lessons, then run npm run generate:books.
% canonical-source-hash: fnv1a64:09003c6e89e4b1d4
% canonical-lessons: JA-C442-kekkan, JA-C442-kinniku, JA-C442-ketsueki, JA-C442-sebone, JA-C442-kansetsu

\chapter{Inside the Body}
\label{ch:ja-inside-the-body}
\hlchaptermodality{442}

\begin{chapteropening}
\textbf{By the end of this chapter:} \emph{I can name the muscles, the blood and the joints.}

Complete the last lesson of chapter 442: I can name the muscles, the blood and the joints.
\end{chapteropening}

\section[\texorpdfstring{\ja{けっかん} (kekkan)}{けっかん (kekkan)}]{\ja{けっかん} (kekkan) — a blood vessel}

\label{lesson:JA-C442-kekkan}



\begin{quote}
Before the new one: say the Japanese for a lazy person, then the Japanese for to fool around.
\end{quote}

\subsection*{You'll want to know: \ja{けっかん}}
\textbf{\ja{けっかん}} — \emph{kekkan} — \textquotedblleft{}a blood vessel\textquotedblright{}.

\begin{quote}
\textbf{\ja{いしゃは} \ja{あしの} \ja{けっかんを} \ja{しらべました。}} — \emph{isha wa ashi no kekkan o shirabemashita} — \textquotedblleft{}The doctor examined the blood vessels in my leg.\textquotedblright{}
\end{quote}

\begin{cousinweb}
\textbf{\ja{けっかん}} (\emph{kekkan}) is Sino-Japanese (kango): two Chinese characters read with their on'yomi, readings borrowed from Chinese. They are \emph{ke}, \textquotedblleft{}blood\textquotedblright{} (Middle Chinese \emph{xwet}) and \emph{kan}, \textquotedblleft{}tube\textquotedblright{} (\emph{kwanX}) (KANJIDIC2; Middle Chinese in Baxter's transcription).
\end{cousinweb}

\begin{grammarlens}[title={What you've built}]
One more word for describing what you lived and felt.
\end{grammarlens}

\subsection*{Your turn}
\begin{itemize}
\raggedright
  \item \emph{Say it:} \emph{kekkan}
  \item \emph{Say it:} \emph{kekkan}, once more
  \item \emph{Say it:} \emph{isha wa ashi no kekkan o shirabemashita} — The doctor examined the blood vessels in my leg.
  \item \emph{Recall:} say the Japanese for a lazy person, then the Japanese for to fool around, then say \emph{kekkan} again
\end{itemize}

\begin{tcolorbox}[breakable,colback=teal!4,colframe=teal!35!black,title={Before you move on}]
What does \ja{けっかん} mean? (\textquotedblleft{}A blood vessel\textquotedblright{}.) Say it once more.
\end{tcolorbox}
\section[\texorpdfstring{\ja{きんにく} (kinniku)}{きんにく (kinniku)}]{\ja{きんにく} (kinniku) — a muscle}

\label{lesson:JA-C442-kinniku}



\begin{quote}
Before the new one: say the Japanese for to fool around, then the Japanese for a blood vessel.
\end{quote}

\subsection*{You'll want to know: \ja{きんにく}}
\textbf{\ja{きんにく}} — \emph{kinniku} — \textquotedblleft{}a muscle\textquotedblright{}.

\begin{quote}
\textbf{\ja{まいにち} \ja{はしって、あしの} \ja{きんにくが} \ja{つよく} \ja{なりました。}} — \emph{mainichi hashitte, ashi no kinniku ga tsuyoku narimashita} — \textquotedblleft{}I ran every day, and my leg muscles got stronger.\textquotedblright{}
\end{quote}

\begin{cousinweb}
\textbf{\ja{きんにく}} (\emph{kinniku}) is Sino-Japanese (kango): two Chinese characters read with their on'yomi, readings borrowed from Chinese. They are \emph{kin}, \textquotedblleft{}muscle\textquotedblright{} (Middle Chinese \emph{kj+n}) and \emph{niku}, \textquotedblleft{}meat\textquotedblright{} (\emph{nyuwk}) (KANJIDIC2; Middle Chinese in Baxter's transcription).
\end{cousinweb}

\begin{grammarlens}[title={What you've built}]
One more word for describing what you lived and felt.
\end{grammarlens}

\subsection*{Your turn}
\begin{itemize}
\raggedright
  \item \emph{Say it:} \emph{kinniku}
  \item \emph{Say it:} \emph{kinniku}, once more
  \item \emph{Say it:} \emph{mainichi hashitte, ashi no kinniku ga tsuyoku narimashita} — I ran every day, and my leg muscles got stronger.
  \item \emph{Recall:} say the Japanese for to fool around, then the Japanese for a blood vessel, then say \emph{kinniku} again
\end{itemize}

\begin{tcolorbox}[breakable,colback=teal!4,colframe=teal!35!black,title={Before you move on}]
What does \ja{きんにく} mean? (\textquotedblleft{}A muscle\textquotedblright{}.) Say it once more.
\end{tcolorbox}
\section[\texorpdfstring{\ja{けつえき} (ketsueki)}{けつえき (ketsueki)}]{\ja{けつえき} (ketsueki) — blood (in medicine)}

\label{lesson:JA-C442-ketsueki}



\begin{quote}
Before the new one: say the Japanese for a blood vessel, then the Japanese for a muscle.
\end{quote}

\subsection*{You'll want to know: \ja{けつえき}}
\textbf{\ja{けつえき}} — \emph{ketsueki} — \textquotedblleft{}blood (in medicine)\textquotedblright{}.

\begin{quote}
\textbf{\ja{びょういんで} \ja{けつえきを} \ja{しらべました。}} — \emph{byōin de ketsueki o shirabemashita} — \textquotedblleft{}They tested my blood at the hospital.\textquotedblright{}
\end{quote}

\begin{cousinweb}
\textbf{\ja{けつえき}} (\emph{ketsueki}) is Sino-Japanese (kango): two Chinese characters read with their on'yomi, readings borrowed from Chinese. They are \emph{ketsu}, \textquotedblleft{}blood\textquotedblright{} (Middle Chinese \emph{xwet}) and \emph{eki}, \textquotedblleft{}fluid\textquotedblright{} (\emph{yek}) (KANJIDIC2; Middle Chinese in Baxter's transcription).
\end{cousinweb}

\begin{grammarlens}[title={What you've built}]
One more word for describing what you lived and felt.
\end{grammarlens}

\subsection*{Your turn}
\begin{itemize}
\raggedright
  \item \emph{Say it:} \emph{ketsueki}
  \item \emph{Say it:} \emph{ketsueki}, once more
  \item \emph{Say it:} \emph{byōin de ketsueki o shirabemashita} — They tested my blood at the hospital.
  \item \emph{Recall:} say the Japanese for a blood vessel, then the Japanese for a muscle, then say \emph{ketsueki} again
\end{itemize}

\begin{tcolorbox}[breakable,colback=teal!4,colframe=teal!35!black,title={Before you move on}]
What does \ja{けつえき} mean? (\textquotedblleft{}Blood (in medicine)\textquotedblright{}.) Say it once more.
\end{tcolorbox}
\section[\texorpdfstring{\ja{せぼね} (sebone)}{せぼね (sebone)}]{\ja{せぼね} (sebone) — the spine; the backbone}

\label{lesson:JA-C442-sebone}



\begin{quote}
Before the new one: say the Japanese for a muscle, then the Japanese for blood.
\end{quote}

\subsection*{You'll want to know: \ja{せぼね}}
\textbf{\ja{せぼね}} — \emph{sebone} — \textquotedblleft{}the spine; the backbone\textquotedblright{}.

\begin{quote}
\textbf{\ja{そふの} \ja{せぼねは} \ja{まがっています。}} — \emph{sofu no sebone wa magatte imasu} — \textquotedblleft{}My grandfather's spine is bent.\textquotedblright{}
\end{quote}

\begin{cousinweb}
\textbf{\ja{せぼね}} (\emph{sebone}) is a native Japanese compound (wago), inherited rather than borrowed from Chinese: \emph{se}, \textquotedblleft{}back\textquotedblright{} and \emph{bone}, \textquotedblleft{}bone (hone)\textquotedblright{} (JMdict, KANJIDIC2).
\end{cousinweb}

\begin{grammarlens}[title={What you've built}]
One more word for describing what you lived and felt.
\end{grammarlens}

\subsection*{Your turn}
\begin{itemize}
\raggedright
  \item \emph{Say it:} \emph{sebone}
  \item \emph{Say it:} \emph{sebone}, once more
  \item \emph{Say it:} \emph{sofu no sebone wa magatte imasu} — My grandfather's spine is bent.
  \item \emph{Recall:} say the Japanese for a muscle, then the Japanese for blood, then say \emph{sebone} again
\end{itemize}

\begin{tcolorbox}[breakable,colback=teal!4,colframe=teal!35!black,title={Before you move on}]
What does \ja{せぼね} mean? (\textquotedblleft{}The spine; the backbone\textquotedblright{}.) Say it once more.
\end{tcolorbox}
\section[\texorpdfstring{\ja{かんせつ} (kansetsu)}{かんせつ (kansetsu)}]{\ja{かんせつ} (kansetsu) — a joint (of the body)}

\label{lesson:JA-C442-kansetsu}



\begin{quote}
Before the new one: say the Japanese for blood, then the Japanese for the spine.
\end{quote}

\subsection*{You'll want to know: \ja{かんせつ}}
\textbf{\ja{かんせつ}} — \emph{kansetsu} — \textquotedblleft{}a joint (of the body)\textquotedblright{}.

\begin{quote}
\textbf{\ja{さむい} \ja{あさは} \ja{ひざの} \ja{かんせつが} \ja{いたみます。}} — \emph{samui asa wa hiza no kansetsu ga itamimasu} — \textquotedblleft{}On cold mornings my knee joints ache.\textquotedblright{}
\end{quote}

\begin{cousinweb}
\textbf{\ja{かんせつ}} (\emph{kansetsu}) is Sino-Japanese (kango): two Chinese characters read with their on'yomi, readings borrowed from Chinese. They are \emph{kan}, \textquotedblleft{}connection\textquotedblright{} (Middle Chinese \emph{kwaen}) and \emph{setsu}, \textquotedblleft{}joint\textquotedblright{} (\emph{tset}) (KANJIDIC2; Middle Chinese in Baxter's transcription).
\end{cousinweb}

\begin{grammarlens}[title={What you've built}]
One more word for describing what you lived and felt.
\end{grammarlens}

\subsection*{Your turn}
\begin{itemize}
\raggedright
  \item \emph{Say it:} \emph{kansetsu}
  \item \emph{Say it:} \emph{kansetsu}, once more
  \item \emph{Say it:} \emph{samui asa wa hiza no kansetsu ga itamimasu} — On cold mornings my knee joints ache.
  \item \emph{Recall:} say the Japanese for blood, then the Japanese for the spine, then say \emph{kansetsu} again
\end{itemize}

\begin{tcolorbox}[breakable,colback=teal!4,colframe=teal!35!black,title={Before you move on}]
What does \ja{かんせつ} mean? (\textquotedblleft{}A joint (of the body)\textquotedblright{}.) Say it once more.
\end{tcolorbox}
